\section{Limitations and Threats to Validity}\label{sec:limitations}

\paragraph{Pilot is under-powered by design.} The v0.3 pilot was scoped to $n{=}5$/domain ($N_{\text{paired}}{=}20$), $K{=}3$, max\_tokens$=200$ to fit a \$20 Haiku budget. With $n{=}5$ the exact sign-flip permutation test has a hard floor of $p{=}0.0312$ and Holm--Bonferroni with $m{=}4$ floors the smallest possible adjusted $p$ at $0.125$. \emph{No primary hypothesis (H1--H4) can cross the pre-registered Holm $p<0.05$ bar at this sample size}, regardless of effect direction or magnitude. A properly powered run at $n{\sim}20$/domain is required to claim or rule out a stable architectural effect.

\paragraph{The v0.3 pilot reports a directional-negative outcome.} On all four primary domains, the \texttt{haiku\_cascade} arm \emph{trailed} the clean-substrate \texttt{haiku\_bare} arm (Hedges' $g \in [-0.90, -0.22]$). The H5.v3 random-effects pool came out at $g=-0.46$ (95\% CI $[-0.93, +0.01]$); H6.v3 and H7.v3 also trailed across all four domains. We interpret this honestly:
\begin{itemize}
  \item The v0.2 ``cascade beats bare'' lift, on a clean Haiku CLI substrate that no longer leaks Claude-Code context into the bare control, does not reproduce. This is consistent with the v0.2 adversarial review's specific prediction.
  \item Against budget-matched (H6.v3) and protocol-matched (H7.v3) controls, the cascade is not preferred on this small sample. Both ``+K compute on a single pass'' and ``a generic 2-pass revise'' are simpler interventions that, at this $n$, perform at least as well as the cascade.
  \item The within-cascade H8.v3 test \emph{is} directionally positive ($\Delta=+0.004$, $g=+0.21$) on the $n=3$ items where event-gated commit chose to revise --- but $n=3$ is far too small to adjudicate causality.
\end{itemize}
Several specific ablations remain to disambiguate ``the cascade hurts'' from ``the pilot is too small / the operators are mistuned'': (a) does an $n{=}20$/domain replication shrink the negative effect or sharpen it? (b) does raising the event-gating $\Delta F$ threshold (so revision fires on more items) recover the H8.v3 signal at meaningful $n$? (c) does an aspect-prior or active-storehouse warm restart on the failing domains improve the cascade selectively? These are the v0.4 falsification ablations, not arguments against this paper's negative finding.

\paragraph{Cascade $K$ deferred from $K{=}4$ to $K{=}3$.} The pilot used $K{=}3$ to keep total Haiku spend under cap. v0.1 ran at $K{=}4$. The v0.3 run is scoped at $K{=}3$; a follow-up at $K{=}4$ on a properly powered $n{=}20$/domain pull will eliminate this confound.

\paragraph{Scoring functions are local proxies, not gold standards.} We deliberately use local deterministic scorers to keep the contrast fair across arms. They are, however, proxies: POEMetric is a more nuanced rubric than what our \texttt{score\_poetry\_gen} function captures, and the Wittgenstein aspect-multiplicity score is constructed from cosine to author-supplied aspect targets rather than from human raters. A follow-up with human raters (or a frontier-model judge) is necessary before the absolute scores can be interpreted as POEMetric scores in the original sense; our \emph{contrasts} between arms remain valid because the same scorer is applied uniformly to all arms.

\paragraph{Prompt-engineering confound.} The Claude Code CLI surrounds every prompt with a Claude Code system context (visible in the cache-creation token count). This system context could in principle help or hurt the Haiku arm in domain-specific ways. We held the user prompt constant across arms and used Claude Code's own \texttt{-p} non-interactive mode (the most LLM-API-like surface), but we cannot prove the system prompt does not bias the Haiku arm's outputs.

\paragraph{Model substrate confound.} Qwen2-1.5B-Instruct is the only local LLM we tested. The PCE design is substrate-agnostic, but a wider sweep (Qwen2-7B, Llama-3-8B, Mistral-Nemo) is needed to confirm the cascade contributes consistently across substrate scales.

\paragraph{Internal-validity vs external-validity.} The H6 test (within-PCE event vs no-event) is a within-condition test and does not establish that the recursive layer would help under a different scoring function or a different prompt distribution. We report it as an \emph{internal} validity check, not as a generalisable claim.

\paragraph{Single-author authorship of the benchmark items.} The Wittgenstein aspect-shift items are author-constructed. We avoided cherry-picking by stratifying across canonical poetic devices (metaphor, synecdoche, allegory, etc.), but the items have not been reviewed by independent literary scholars. We publish them in Appendix~\ref{app:benchmark_items} so the community can replace or augment them.
